\subsection{nml\_output}
\label{SEC:nml_output}
\inp*[nmlOutput]{}[nml\_output][]

This \hyperref[SEC:User_Input]{namelist block}
is used to control various output files.

\begin{nmlblocklst}[
,label={LST:nml_output}
,caption={Namelist nml\_output.}
]
&nml_output
  Time_Scaling_Factor                =  1.0
  Display_Interval                   =  1
  Residual_Interval                  =  0
  Restart_Interval                   =  0
  Restart_Frequency                  = -1.0  ! tstar
  Restart_Output_Uses_Host_Roots     =  T
  Restart_Output_Disabled            =  F
  Output_Dir                         = "."
  Output_Interval                    = -5
  Output_Frequency                   = -1.0  ! tstar
  Output_Order                       = -3
  Output_Summary_File                =  F
  Dmat_Grad_Temp_Order               = -1
  Loc_Output_Pts                     =  0
  CGNS_High_Order_Output             =  F
  CGNS_Mixed_Output                  =  T
  CGNS_Output_Disabled               =  F
  Probes_Collection_Interval         =  0
  Probes_Write_Interval              =  100
  Point_Probes(1)%xyz(1:3)           =  [-999.,-999.,-999.] ! grid units
  Rake_Probes(1)%num_probes          =  0
  Rake_Probes(1)%xyz_beg(1:3)        =  [-999.,-999.,-999.] ! grid units
  Rake_Probes(1)%xyz_end(1:3)        =  [-999.,-999.,-999.] ! grid units
  Cell_Probes(1)%xyz(1:3)            =  [-999.,-999.,-999.] ! grid units
  Profile_IO_All                     =  F
  Profile_IO_Restart                 =  F
  Profile_IO_CGNS                    =  F
  Profile_IO_CGNS_Subs               =  F
  Lustre_Stripe_Count                =  64
  Lustre_Stripe_Size                 = "4m"
/
\end{nmlblocklst}

%==========================================================================

\begin{nmldescitem}
Time_Scaling_Factor = 1.0  ! [1/s]
\end{nmldescitem}
\inp*[nmlOutput::TimeScalingFactor]{}[Time\_Scaling\_Factor][]

\begin{nmldescbody}
This scaling factor is used to define a non-dimensional time
(\key*[key:tstar]{\tstar}[t *@\tstar])
\begin{gather}
  \label{EQN:tstar}
  \tstar = ( t - t_0 ) \cdot \inptext{Time\_Scaling\_Factor}
\end{gather}
where $t$ is the dimensional time (in seconds).
%
If time-averaging has been activated, then $t_0$ is the dimensional time
(in seconds) when time-averaging began.
Otherwise, $t_0=0$.
%
This non-dimensional time (\tstar)
is what is written to standard output and the residual file.
It is also used by the following input variables to
determine when certain events should occur:
\begin{itemize}
\item
  \inp[nmlPhysics::MovingRefFrame/Spinup/Tstar]{Moving\_Ref\_Frame\%Spinup\%Tstar} for the moving reference frame startup,
\item
  \inp[nmlTimeAve::StartAveragingTstar]{Start\_Averaging\_Tstar}
  for time averaging startup,
\item
  \inp[nmlOutput::RestartFrequency]{Restart\_Frequency} of restart files,
\item
  \inp[nmlOutput::OutputFrequency]{Output\_Frequency} of CGNS files,
  and
\item
  \inp[nmlTimestep::FinalTime]{Final\_Time} for when to stop a simulation.
\end{itemize}


For example, many simulations are based on a number of flow-through periods.
%
In this case, the \inptext{Time\_Scaling\_Factor} should be set to
the value of $U/L$, where
  $U$ is the convective velocity and
  $L$ is the flow-through length of the domain.
\end{nmldescbody}

%==========================================================================

\begin{nmldescitem}
Display_Interval = 1
\end{nmldescitem}
\inp*[nmlOutput::DisplayInterval]{}[Display\_Interval][]

\begin{nmldescbody}
Number of time steps between writing the time-stepping, residual,
and convergence statistics to standard output.
%
The input is treated as follows:
\\
\begin{tabular}{ @{\qquad} r @{ \quad } p{0.75\linewidth} @{} }
  $< 0$ & Reset to
          $\inp[nmlTimestep::NumTimesteps]{Num\_Timesteps} /
           | \inptext{Display\_Interval} |$.\\
  $= 0$ & Reset to 1.\\
  $> 0$ & Use as-is.
\end{tabular}

\NOTE
When running the Taylor-Green vortex problem or the
Taylor-Couette problem, this input is also used as the interval for
writing integrated kinetic energy dissipation rate (KEDR)
statistics to the file ``TaylorGreen\_KEDR.dat'' or
``TaylorCouette\_KEDR.dat'', respectively.
\end{nmldescbody}

%==========================================================================

\begin{nmldescitem}
Residual_Interval = 0
\end{nmldescitem}
\inp*[nmlOutput::ResidualInterval]{}[Residual\_Interval][]

\begin{nmldescbody}
Number of time steps between writing the residual and convergence
statistics for all conserved and primitive variables to the
residual file ``residual.dat''.
%
The input is treated as follows:
\\
\begin{tabular}{ @{\qquad} r @{ \quad } p{0.75\linewidth} @{} }
  $< 0$ & Reset to
          $\inp[nmlTimestep::NumTimesteps]{Num\_Timesteps} /
           | \inptext{Residual\_Interval} |$.\\
  $= 0$ & Reset to
          $\inp[nmlTimestep::NumTimesteps]{Num\_Timesteps} / 500$.\\
  $> 0$ & Use as-is.
\end{tabular}
\end{nmldescbody}

%==========================================================================

\begin{nmldescitem}
Restart_Interval = 0
\end{nmldescitem}
\inp*[nmlOutput::RestartInterval]{}[Restart\_Interval][]

\begin{nmldescbody}
Specifies the number of time steps between writing restart files.
%
Note that a restart file is always written before the first (non-restarted)
time step and after completing the last time step.
%
With the default value of 0, no intermediate files will be written,
only the restarts before the first time step and after the final time step.
%
After successfully writing a new restart file, a
symbolic link will be created that points to the latest restart
file.
%
The name of this link will be:
\newline\indent
``\inp[nmlOutput::OutputDir]{Output\_Dir}\slash{Restart\_files}\slash{out.current.rst}'' 
\end{nmldescbody}

%==========================================================================

\begin{nmldescitem}
Restart_Frequency = -1.0  ! tstar
\end{nmldescitem}
\inp*[nmlOutput::RestartFrequency]{}[Restart\_Frequency][]

\begin{nmldescbody}
If $>0$, this value specifies the non-dimensional time
(\key[key:tstar]{\tstar}[t *@\tstar]) interval, \eref{EQN:tstar},
to write restart files.
%
This is complimentary to \inptext{Restart\_Interval}
and will write additional restart files.
%
After successfully writing a new restart file, a
symbolic link will be created that points to the latest restart
file.
%
The name of this link will be:
\newline\indent
``\inp[nmlOutput::OutputDir]{Output\_Dir}\slash{Restart\_files}\slash{out.current.rst}'' 
\end{nmldescbody}

%==========================================================================

\begin{nmldescitem}
Restart_Output_Uses_Host_Roots = T
\end{nmldescitem}
\inp*[nmlOutput::RestartOutputUsesHostRoots]{}[Restart\_Output\_Uses\_Host\_Roots][]

\begin{nmldescbody}
Indicates which MPI processes are involved in writing to the
restart file.
%
If true, then only the root processes on each host node will
preform the output.
%
Otherwise, all MPI processes will be involved in the output.

\NOTE
This is only applicable for parallel simulations run across
multiple physical computers, \forexample multiple nodes on
the NASA HECC supercomputers.
%
Only change this if you encounter MPI errors that occur when
trying to write restart files.
%
Changing this to false will most likely result in higher
I\slash O file contention and possibly increased memory usage.
\end{nmldescbody}

%==========================================================================

\begin{nmldescitem}
Restart_Output_Disabled = F
\end{nmldescitem}
\inp*[nmlOutput::RestartOutputDisabled]{}[Restart\_Output\_Disabled][]

\begin{nmldescbody}
If set to true, this option will prevent any instantaneous or time-averaged
restart files from being written.
%
This is intended for simulations that will complete with the allotted
resources and will not need to be restarted later.
\end{nmldescbody}

%==========================================================================

\begin{nmldescitem}
Output_Dir = "."
\end{nmldescitem}
\inp*[nmlOutput::OutputDir]{}[Output\_Dir][]

\begin{nmldescbody}
File path to the directory used for writing output files.

\NOTE
If one is available, it is recommended that the output directory be
located on a Lustre file system in order to get the best I\slash O
performance possible.
\end{nmldescbody}

%==========================================================================

\begin{nmldescitem}
Output_Interval = -5
\end{nmldescitem}
\inp*[nmlOutput::OutputInterval]{}[Output\_Interval][]

\begin{nmldescbody}
This controls the number of time steps between writing CGNS solution files.
The input value is interpreted as follows:
\\
\begin{tabular}{ @{\qquad} r @{ \quad } p{0.85\linewidth} @{} }
  $<0$ & For most situations, every
         $\inp[nmlTimestep::NumTimesteps]{Num\_Timesteps} /
           | \inptext{Output\_Interval} |$
         time steps.\\
  $=0$ & For most situations, every
         $\inp[nmlTimestep::NumTimesteps]{Num\_Timesteps}/50$ time steps.\\
  $>0$ & For all situations, every
         \inptext{Output\_Interval} time steps.
\end{tabular}
\\
The first two options result in a fixed number of output files
each time the code is run, assuming the code does not terminate
early.


\NOTE
For the following special cases, the input value is interpreted differently:
\begin{itemize}
%---------------------------------------------------------------------
\item 
If either
  $\inp[nmlOutput::OutputFrequency]{Output\_Frequency}>0$ or
  $\inp[nmlOutput::CGNSOutputDisabled]{CGNS\_Output\_Disabled}=\texttt{T}$,
the action associated with those inputs will take precedence
and this input will not be used.
%---------------------------------------------------------------------
\item
If running any of the transport problems (isentropic vortex,
density\slash entropy advection), this input value is used to
determine when a CGNS solution file will be written as follows:
\\
\begin{tabular}{ @{\qquad} r @{ \quad } p{0.75\linewidth} @{} }
  $<0$ & At the completion of every $|\inptext{Output\_Interval}|$
         flow-through periods.\\
  $=0$ & At the completion of every flow-through period.
\end{tabular}
%---------------------------------------------------------------------
\item
If running either the Taylor-Green vortex problem or the
Taylor-Couette problem, CGNS solution files will automatically be
written at certain specific values of
non-dimensional time (\key[key:tstar]{\tstar})
that are of interest
to each specific problem. If $\inptext{Output\_Interval}=0$,
only these problem specific CGNS solution files will be written;
otherwise ($\inptext{Output\_interval}\neq 0$) this behaves the
same as ``most situations'' in addition to writing the
problem-specific CGNS solution files.
%---------------------------------------------------------------------
\end{itemize}

\end{nmldescbody}

%==========================================================================

\begin{nmldescitem}
Output_Frequency = -1.0  ! tstar
\end{nmldescitem}
\inp*[nmlOutput::OutputFrequency]{}[Output\_Frequency][]

\begin{nmldescbody}
If $>0$, this value specifies the non-dimensional time
(\key[key:tstar]{\tstar}[t *@\tstar]) interval, \eref{EQN:tstar},
to write CGNS files.
%
\inp[nmlOutput::OutputInterval]{Output\_Interval} will be disabled.
%% Users need not know this detail:
%  by setting it to $\inp[nmlTimestep::NumTimesteps]{Num\_Timesteps}+1$.
%
\end{nmldescbody}

%==========================================================================

\begin{nmldescitem}
Output_Order = -3
\end{nmldescitem}
\inp*[nmlOutput::OutputOrder]{}[Output\_Order][]

\begin{nmldescbody}
Specifies the degree of the polynomial space used to represent the solution
in the CGNS solution files. This allows the ability to over-sample the
solution polynomial within each cell which provides a smoother solution 
when visualizing within your favorite post-processing software.

\NOTE
If using \inp[nmlOutput::CGNSHighOrderOutput]{CGNS\_High\_Order\_Output}, the
output order will be limited to a maximum of 4, since that is the
maximum order supported by the CGNS library.

\NOTE
If
$\inptext{Output\_Order} < \inp[nmlNumerics::SolutionOrder]{Solution\_Order}$,
the solution order will be used instead.
\end{nmldescbody}

%==========================================================================

\begin{nmldescitem}
Output_Summary_File = F
\end{nmldescitem}
\inp*[nmlOutput::OutputSummaryFile]{}[Output\_Summary\_File][]

\begin{nmldescbody}
Setting this to true will create the file ``summary.log'' when the
simulation successfully stops.
%
This file contains summarized information about the simulation like:
\begin{itemize}[noitemsep,nolistsep]
\item
  computational cost
\item
  solution order
\item
  number of cells
\item
  degrees of freedom
\item
  total time or \key[key:tstar]{\tstar}[t *@\tstar]
  elapsed during this simulation
\item
  final time or \key[key:tstar]{\tstar}[t *@\tstar]
  at the end of the simulation\\
  (\thatis combining the elapsed value with that from the restart file)
\item
  the current length of time averaging (if applicable)
\item
  and more
\end{itemize}

When running the channel flow test case using the internal grid generator,
it will also contain detailed grid information such as the number of
grid cells and degrees of freedom in each direction, as well as their
respective spacings.
\end{nmldescbody}

%==========================================================================

\begin{nmldescitem}
Dmat_Grad_Temp_Order = -1
\end{nmldescitem}
\inp*[nmlOutput::DmatGradTempOrder]{}[Dmat\_Grad\_Temp\_Order][]

\begin{nmldescbody}
If outputting the derivative of temperature, this is the quadrature order
used for computing temperature before projecting down to the solution order
and computing the derivative.

\NOTE
If $\inptext{Dmat\_Grad\_Temp\_Order} < 0$,
\inp[nmlNumerics::SolutionOrder]{Solution\_Order}
will be used instead.
\end{nmldescbody}

%==========================================================================

\begin{nmldescitem}
Loc_Output_Pts = 0
\end{nmldescitem}
\inp*[nmlOutput::LocOutputPts]{}[Loc\_Output\_Pts][]

\begin{nmldescbody}
Specifies the location of the \key{output nodal points}[points!output]
 within the grid cells when writing the CGNS solution files.
\\
\begin{tabular}{ @{\qquad} r @{ = } p{0.85\linewidth} @{} }
  0 & \key{Equidistant points}[points!equidistant] nodes \\
  1 & \key{Legendre-Gauss (Gauss) points}[points!Legendre-Gauss (Gauss)] nodes \\
  2 & \key{Legendre-Gauss-Lobatto (Lobatto) points}[points!Legendre-Gauss-Lobatto (Lobatto)] nodes
\end{tabular}

\NOTE
Using Legendre-Gauss nodes for the CGNS solution output files
will result in empty space between all cells because
those points do not extend all the way to the cell edge and
there is no connectivity information between cells.
%
This only affects your ability to visualize the solution
and does not affect the solution accuracy in any way.
\end{nmldescbody}

%==========================================================================

%\begin{nmldescitem}
%Continuous_Output = F
%\end{nmldescitem}
%\inp*[nmlOutput::ContinuousOutput]{}[Continuous\_Output][]
%
%\begin{nmldescbody}
%Specifies whether the output routine should attempt to create a smooth,
%continuous solution across all cells before writing to the
%CGNS solution file.
% 
%This option is intended for use with
%\inp[nmlOutput::LocOutputPts]{Loc\_Output\_Pts}=1 (Legendre-Gauss nodes),
%which have empty space between all cells because
%those points do not extend all the way to the cell edge and
%there is no connectivity information between cells.
%%
%When activated, additional output points will be added within
%each cell to fill this gap.
%%
%This will create CGNS solution files that are significantly larger in file
%size compared to the regular CGNS solution files.
%
%\NOTE
%This only works with a single MPI process.
%%
%So far, tests of this feature have been underwhelming with regards
%to creating a cleaner visualization of the solution for things
%like isosurfaces.
%
%\end{nmldescbody}

%==========================================================================

%\begin{nmldescitem}
%Interpolate_Before_Output_Variable = T
%\end{nmldescitem}
%\inp*[nmlOutput::InterpolateBeforeOutputVariable]{}[Interpolate\_Before\_Output\_Variable][]
%
%\begin{nmldescbody}
%Determines whether interpolation to the output points is done
%before or after the output variables are computed.
%
%\begin{tabular}{ @{} r @{ = } p{0.85\linewidth} @{} }
%  T & Interpolation is done before the output variable is computed.
%      Flow variables are interpolated from the solution points
%      to the output points, then output quantities are computed
%      at the output points.
%      This provides consistency between the various output variables
%      at the output points.\\
%  F & Interpolation is done after the output variable is computed.
%      Output quantities are computed from the flow variables
%      at the solution points, then interpolated to the output points.
%\end{tabular}
%
%In GFR, the conserved variables are represented as polynomials
%within each element.
%%
%However, derived quantities such as velocity, temperature, and
%pressure cannot be exactly represented as a polynomial
%because they involve nonlinear relations or fractions of polynomials.
%%
%If one interpolates the conserved variables to the output points
%and then computes the derived quantities, all the output variables
%at the output points will be self-consistent.
%%
%If one instead computes the derived quantities at the solution points
%and uses those approximate polynomials to interpolate onto the
%output points, the conserved and derived values at the output
%points will not be self-consistent.
%
%\end{nmldescbody}

%==========================================================================

\begin{nmldescitem}
CGNS_High_Order_Output = F
\end{nmldescitem}
\inp*[nmlOutput::CGNSHighOrderOutput]{}[CGNS\_High\_Order\_Output][]

\begin{nmldescbody}
Setting this to true will write the CGNS grid file (solution.grid.cgns)
using high-order cells. If not enabled (the default), the CGNS grid file
breaks each cell into linear sub-cells.

\NOTE
Visualizing high-order CGNS solution files requires specialized software,
such as Tecplot 2023 or newer.
\end{nmldescbody}

%==========================================================================

\begin{nmldescitem}
CGNS_Mixed_Output = T
\end{nmldescitem}
\inp*[nmlOutput::CGNSMixedOutput]{}[CGNS\_Mixed\_Output][]

\begin{nmldescbody}
This determines how cell connectivity is handled in high-order CGNS grid files.
%
If true, all grid cells will be combined into a single zone
regardless of cell geometry or order.
%
If false, each individual combination of cell geometry and order
will be written to its own zone.
%
Some post-processing software may prefer the latter format.
\end{nmldescbody}

%==========================================================================

\begin{nmldescitem}
CGNS_Output_Disabled = F
\end{nmldescitem}
\inp*[nmlOutput::CGNSOutputDisabled]{}[CGNS\_Output\_Disabled][]

\begin{nmldescbody}
Completely disables CGNS solutions from being output at the output interval.

On some systems with a large number of processors,
the output of CGNS files can be slow.
%
It can be advantageous to instead output restart files,
and later use
\inp[nmlUtility::ConvertRestartToCGNS]{Convert\_Restart\_To\_CGNS}
in the \inptext{nml\_utility} namelist block of \sref{SEC:nml_utility}
to have GFR batch convert those restart files to CGNS.
%
This post-processing step can be done using a minimal number
of processors, because most of the work is simple file I\slash O.

\end{nmldescbody}

%==========================================================================

\begin{nmldescitem}
Probes_Collection_Interval = 0
\end{nmldescitem}
\inp*[nmlOutput::ProbesCollectionInterval]{}[Probes\_Collection\_Interval][]

\begin{nmldescbody}
If $>0$, this value specifies the iteration interval used to collect
time-history data for each of the probes. The information collected at this
interval is saved to a buffer and only written out after completing the number
of collections specified by the input
\inp[nmlOutput::ProbesWriteInterval]{Probes\_Write\_Interval}.
%
With the default value of 0, no time-history probe data will be collected or
output.
\end{nmldescbody}

%==========================================================================

\begin{nmldescitem}
Probes_Write_Interval = 100
\end{nmldescitem}
\inp*[nmlOutput::ProbesWriteInterval]{}[Probes\_Write\_Interval][]

\begin{nmldescbody}
This value specifies how often to write out the buffer used to collect
time-history data for the probes.

\NOTE
This interval is relative to the number of collections completed, and is not
relative to the simulation iteration.
\end{nmldescbody}

%==========================================================================

\begin{nmldescitem}
Point_Probes(1)%xyz(1:3) = [-999.,-999.,-999.]  ! input grid units
\end{nmldescitem}
\inp*[nmlOutput::PointProbes]{}[Point\_Probes][]

\begin{nmldescbody}
This specifies the location of a single point \key[key:probes]{probe}[probes] to
collect time-history data.
%
The values specified for this input are assumed to be in the units of the input
mesh file (\thatis \emph{grid units} as discussed in \sref{SEC:Grid_Scaling}).
%
Up to 200 point probes may be used.
%
Probes set to the default location are ignored.
%
For more information on the point probe files, see
\sref{SEC:Files:Probes_Output}.

\NOTE
The solution point closest to these coordinates will be used as the source for
the probe data.

\NOTE
If an individual solution point is identified as the source of a point probe by
more than one point\slash rake\slash cell probe, the duplicate probes will be
removed and only a single probe will be collected for that solution point.
%
This will make it seem like there are missing probe files.
\end{nmldescbody}

%==========================================================================

\begin{nmldescitem}
Rake_Probes(1)%Num_Probes   = 0
Rake_Probes(1)%xyz_beg(1:3) = [-999.,-999.,-999.]  ! input grid units
Rake_Probes(1)%xyz_end(1:3) = [-999.,-999.,-999.]  ! input grid units
\end{nmldescitem}
\inp*[nmlOutput::RakeProbes]{}[Rake\_Probes][]

\begin{nmldescbody}
These inputs provide the details necessary to create a rake of point
\key[key:probes]{probes}[probes].
%
The number of probes to create within rake $i$ is specified using the input
\inptext{Rake\_Probes($i$)\%Num\_Probes}.
%
\inptext{Rake\_Probes($i$)\%xyz\_beg(1:3)} and
\inptext{Rake\_Probes($i$)\%xyz\_end(1:3)}
specify the coordinates for the start and end points of the rake, with both of
these coordinate inputs assumed to be in the units of the mesh file (\thatis
\emph{grid units} as discussed in \sref{SEC:Grid_Scaling}).

In the case that $\inptext{Num\_Probes}=1$, the single probe created by the
rake will be located halfway between the start and end points for the rake.
%
Otherwise, probes will be located at both the start and end points for the
rake, and $\inptext{Num\_Probes}-2$ probes will be equally spaced along the
line between the start and end points.

Up to 10 rake probes may be used, with no limit on the number of probes per rake.
%
Rakes with $\inptext{Num\_Probes}=0$ or $\inptext{xyz\_beg}=\inptext{xyz\_end}$
are ignored, and no time-history probe data will be collected or
output for the rake.
%
For more information on the rake probe files, see
\sref{SEC:Files:Probes_Output}.

\NOTE
For each individual rake probe, the solution point closest to the probe
coordinates will be used as the source for the probe data.

\NOTE
If an individual solution point is identified as the source of a point probe by
more than one point\slash rake\slash cell probe, the duplicate probes will be
removed and only a single probe will be collected for that solution point.
%
This will make it seem like there are missing probe files.

\NOTE
It should be warned that too many probes (1000+) may have a large negative
impact on code and file system performance.
\end{nmldescbody}

%==========================================================================

\begin{nmldescitem}
Cell_Probes(1)%xyz(1:3) = [-999.,-999.,-999.]  ! input grid units
\end{nmldescitem}
\inp*[nmlOutput::CellProbes]{}[Cell\_Probes][]

\begin{nmldescbody}
This input provides the ability to create a set of point probes for every
solution point within a grid cell.
%
For the grid cell that contains the coordinates of this input, individual point
\key[key:probes]{probes}[probes] will be created for every solution point within
that grid cell, resulting in time-history data for the entire cell.
%
The values specified for this input are assumed to be in the units of the input
mesh file (\thatis \emph{grid units} as discussed in \sref{SEC:Grid_Scaling}).
%
Up to 10 cell probes may be used.
%
Probes set to the default location are ignored.
%
For more information on the cell probe files, see
\sref{SEC:Files:Probes_Output}.

\NOTE
If an individual solution point is identified as the source of a point probe by
more than one point\slash rake\slash cell probe, the duplicate probes will be
removed and only a single probe will be collected for that solution point.
%
This will make it seem like there are missing probe files.

\NOTE
It should be warned that too many cell probes may have a large negative impact
on code and file system performance.
%
A single cell probe for a \PO[3] hexahedral cell will result in 64 separate
point probes, and 125 probes for a \PO[4] hexahedral cell.
\end{nmldescbody}

%==========================================================================

\begin{nmldescitem}
Profile_IO_All = F
\end{nmldescitem}
\inp*[nmlOutput::ProfileIOAll]{}[Profile\_IO\_All][]

\begin{nmldescbody}
If true, sets both
  \inp[nmlOutput::ProfileIOCGNS]{Profile\_IO\_CGNS} and
  \inp[nmlOutput::ProfileIORestart]{Profile\_IO\_Restart}
to true.
%
If false, it has no effect on the other input parameters.
\end{nmldescbody}

%==========================================================================

\begin{nmldescitem}
Profile_IO_Restart = F
\end{nmldescitem}
\inp*[nmlOutput::ProfileIORestart]{}[Profile\_IO\_Restart][]

\begin{nmldescbody}
Activates timers to measure how long it takes to write each restart or
time-averaged restart file, and outputs the results to standard output.
\end{nmldescbody}

%==========================================================================

\begin{nmldescitem}
Profile_IO_CGNS = F
\end{nmldescitem}
\inp*[nmlOutput::ProfileIOCGNS]{}[Profile\_IO\_CGNS][]

\begin{nmldescbody}
Activates timers to measure how long it takes to write each CGNS solution file,
and outputs the results to standard output.
\end{nmldescbody}

%==========================================================================

\begin{nmldescitem}
Profile_IO_CGNS_Subs = F
\end{nmldescitem}
\inp*[nmlOutput::ProfileIOCGNSSubs]{}[Profile\_IO\_CGNS\_Subs][]

\begin{nmldescbody}
Activates timers to measure how long it takes to complete each individual call
to a function from the CGNS library, and outputs the results to standard output.
%Outputs timing information for each individual call to a CGNS function.
\end{nmldescbody}

%==========================================================================

\begin{nmldescitem}
Lustre_Stripe_Count = 64
\end{nmldescitem}
\inp*[nmlOutput::LustreStripeCount]{}[Lustre\_Stripe\_Count][]

\begin{nmldescbody}
Specifies the approximate number of stripes to use when writing a CGNS
solution file or a restart file.

\NOTE
This is only used if \inp[nmlOutput::OutputDir]{Output\_Dir}
is on a Lustre file system.
%
The code will take this value and find the closest integer value satisfying:
$\textrm{mod} \left( \textsl{ncpu}
                   , \inptext{Lustre\_Stripe\_Count}
              \right) = 0$
\end{nmldescbody}

%==========================================================================

\begin{nmldescitem}
Lustre_Stripe_Size = "4m"
\end{nmldescitem}
\inp*[nmlOutput::LustreStripeSize]{}[Lustre\_Stripe\_Size][]

\begin{nmldescbody}
Specifies the number of bytes to store on each stripe before moving to the
next stripe.

\NOTE
This is only used if \inp[nmlOutput::OutputDir]{Output\_Dir}
is on a Lustre file system.
%
The value of \linebreak[4]\inptext{Lustre\_Stripe\_Size}
must be a multiple of 65,536 bytes (64 KB), and the suffixes
`k', `m', or `g' can be used to specify units of
KB, MB, or GB, respectively.

\EXAMPLE
$\textrm{``4m''} \Rightarrow 4~\textrm{MB}
                 \Rightarrow 4,194,304~\textrm{bytes}$
\end{nmldescbody}

%==========================================================================

